\chapter{Declared Physics: TFTP and the Sorcerer's Apprentice}
\label{ch:tftp-physics}

\index{TFTP}
\index{Sorcerer's Apprentice}
\index{physics of a channel}
A grammar does not say what happens to a message after it is sent: whether it
queues, how large the queue is, whether the network loses it, duplicates it or
delays it. Those facts are part of what a protocol means. The example is the
Trivial File Transfer Protocol (TFTP) \citep{rfc783,rfc1350} and its best-known
defect, the \emph{Sorcerer's Apprentice} bug, named and fixed in RFC~1123
\citep{rfc1123}.

This is the book's central claim in its smallest form. As a multiparty session type
\citep{honda2008multiparty}, the original and the fixed sender are the same
well-formed protocol. With the physics declared, LeanFM finds a walk in which one
early timer makes every later block travel twice for the original, and proves
that none exists for the fix.

Status: the verdicts are real \texttt{check\_ctl} results, and the protocol and
its physics are checked Lean in \texttt{SessionExample.lean}. The product of roles
and queues is still built by a small script, \texttt{video/tftp\_model.mjs} in
the companion video repository; deriving it from the declarations is future work.

\section{The grammar}

TFTP moves a file in numbered blocks. The sender sends \texttt{DATA}~$n$; the
receiver answers \texttt{ACK}~$n$; the sender then sends \texttt{DATA}~$n{+}1$.
A block shorter than the full size ends the transfer.

In LeanFM's session-type syntax, with the physics beside it:

\begin{lstlisting}[language={}]
protocol Transfer (Sender, Receiver) {
  rec Block {
    d: Data from Sender to Receiver;
    choice at Receiver {
      Ack from Receiver to Sender cause d;
      continue Block;
    } or {
      LastAck from Receiver to Sender cause d;
    }
  }
}

def tftpPhysics : Physics Role :=
  { channels := [ (.sender, .receiver, { udp 2 with duplication := false })
                , (.receiver, .sender, { udp 2 with duplication := false }) ]
    timers := [{ role := .sender }] }
\end{lstlisting}


Without a timer this is a textbook global type: the receiver chooses,
answering \texttt{Ack} for a full block and \texttt{LastAck} for a short one, and
the projected roles never deadlock or receive an unexpected message. But a real
network loses packets, so the sender needs a timer that resends a block when no
\texttt{ACK} arrives in time. That retransmission is the whole engineering content
of TFTP, and the classic reading has nowhere to put it, because it assumes every
message arrives exactly once, in order.

\section{The physics}

\index{UDP}
\index{queue capacity}
\texttt{tftpPhysics} above declares each channel as \texttt{udp} with one
difference written out. Each direction holds at most two packets, delivered in
either order or dropped, and the sender's timer may fire any time a block is
unacknowledged, since a slow answer looks like a lost one. The network is declared
never to duplicate: every duplicate below is made by the sender. The model has
three blocks. The two senders differ in one line, the reaction to a repeated
\texttt{ACK}:

\begin{description}[style=nextline]
  \item[Original (RFC 783).] Every \texttt{ACK}~$n$, even a repeated one, makes
  the sender send \texttt{DATA}~$n{+}1$.
  \item[Fixed (RFC 1123).] A repeated \texttt{ACK} is ignored. A block is sent
  again only when its own timer fires.
\end{description}

The checked machine is the product of sender, receiver and both queues: 201
reachable states and 891 transitions for each sender. Its events are
\texttt{timeout}, \texttt{drop}, \texttt{data} (a new block delivered),
\texttt{dup}$k$ (block $k$ delivered again, and acknowledged again), \texttt{ack}
(the first \texttt{ACK} for the current block), \texttt{dupack} (an \texttt{ACK}
for an earlier block) and \texttt{done}.

\section{The question}

Suppose the timer fires early once, at the start, so \texttt{DATA}~1 is in flight
twice, and afterwards nothing is lost and the timer never fires. Can block~3 still
arrive twice? From the state after that \texttt{timeout}:

\AccessibleEquation{Not the disjunction of timeout and drop holds along some path until dup three.}{\[
  E\,[\,\lnot(\mathit{timeout} \lor \mathit{drop})\ U\ \mathit{dup}_3\,]
\]}

``Some path has no \texttt{timeout} and no \texttt{drop} until block~3 is
delivered a second time.'' For the original sender LeanFM proves it, with the
witness in \cref{tab:tftp-witness}; the queue columns are read from the states on
the path.

\begin{table}[ht]
\centering
\small
\begin{tabular}{@{}rlp{0.42\textwidth}ll@{}}
\toprule
step & event & meaning & DATA queue & ACK queue \\
\midrule
1 & \texttt{timeout} & timer fires early; DATA 1 sent again & 1, 1 & -- \\
2 & \texttt{data} & DATA 1 delivered & 1 & 1 \\
3 & \texttt{dup1} & DATA 1 delivered again; ACK 1 sent again & -- & 1, 1 \\
4 & \texttt{ack} & ACK 1 received; DATA 2 sent & 2 & 1 \\
5 & \texttt{data} & DATA 2 delivered & -- & 1, 2 \\
6 & \texttt{dupack} & ACK 1 received again; DATA 2 sent again & 2 & 2 \\
7 & \texttt{dup2} & DATA 2 delivered again; ACK 2 sent again & -- & 2, 2 \\
8 & \texttt{ack} & ACK 2 received; DATA 3 sent & 3 & 2 \\
9 & \texttt{data} & DATA 3 delivered & -- & 2, 3 \\
10 & \texttt{dupack} & ACK 2 received again; DATA 3 sent again & 3 & 3 \\
11 & \texttt{dup3} & DATA 3 delivered again; ACK 3 sent again & -- & 3, 3 \\
\bottomrule
\end{tabular}
\caption{LeanFM's witness for the original sender: one early timer, then no
loss and no further timeouts, and every later block is sent and delivered
twice.}
\label{tab:tftp-witness}
\end{table}

The pattern repeats every four steps: a doubled block brings a doubled
\texttt{ACK}, whose first copy sends the next block and whose late second copy sends
it again. Nothing fails and every packet arrives, yet the duplication sustains
itself, like the brooms that keep multiplying, doubling the traffic of real
transfers. For the fixed sender the formula is disproved over all 509 reachable
nodes: the repeated \texttt{ACK}~1 is ignored and the duplication dies out.

\section{Liveness needs fairness}

\index{fairness}
Two more questions are asked of both senders.

\AccessibleEquation{A G E F terminal holds for both senders. A F terminal fails for both senders.}{\[
  AG\,EF\,\mathit{terminal}\ \text{holds},\qquad
  AF\,\mathit{terminal}\ \text{fails}.
\]}

From every state some continuation finishes, but not every path does: LeanFM
returns a lasso in which the timer fires forever, each copy dropped by a full
queue, and nothing is delivered. No protocol can force a network to deliver, so a
promise that the transfer \emph{will} finish needs a stated fairness assumption on
the channel, such as ``a packet sent infinitely often is delivered infinitely
often''.

\section{What the classic session-type reading leaves out}

\index{multiparty session type}
\begin{table}[ht]
\centering
\small
\begin{tabular}{@{}p{0.22\textwidth}p{0.32\textwidth}p{0.36\textwidth}@{}}
\toprule
sender & as a session type & with the physics declared \\
\midrule
original (RFC 783) & well-formed lock step & one early timer doubles every later block \\
fixed (RFC 1123) & the same well-formed lock step & no cascade; checked over all 509 nodes \\
\bottomrule
\end{tabular}
\caption{The same grammar, two meanings.}
\label{tab:tftp-contrast}
\end{table}

The grammar is the same in both rows of \cref{tab:tftp-contrast}; only its meaning
changed. With the theory's unbounded, reliable, ordered queues
\citep{honda2008multiparty} a retransmission timer is unnecessary, and the
Sorcerer's Apprentice cannot even be expressed. LeanFM's position follows from the
actor view:

\begin{itemize}
  \item Every choice is an internal choice of one actor, made only from that
        actor's own history: what it sent, and what arrived in its queue
        (\texttt{prev} and \texttt{cause}).
  \item Every queue has a declared capacity, and a declared policy for when it
        is full.
  \item Every channel declares what the network may do: drop, reorder,
        duplicate, delay.
  \item Liveness claims carry their fairness assumptions as named
        requirements.
\end{itemize}

UDP protocols make the physics impossible to ignore, and TFTP is the smallest with
a famous bug that stays invisible until the queues and the early timer are modeled.
